% Source: ME20_v1.html

\documentclass[11pt]{article}
\usepackage[margin=1in]{geometry}
\usepackage{amsmath}
\usepackage{amssymb}
\usepackage{siunitx}
\usepackage{enumitem}

\title{ME20: Kinetic Theory / Ideal Gas}
\author{}
\date{}

\begin{document}
\maketitle

\section*{ME20: Kinetic Theory / Ideal Gas}

\begin{enumerate}[leftmargin=*,label=\textbf{Question \arabic*.}]

\item \hfill \textit{10 pts}

A gasoline engine takes in 8.87 kJ of heat energy and has a thermal efficiency of 0.37. Calculate the value of  $Q_C$  (in kJ). Be careful of the sign.

\textit{Correct answer:} -5.5881

\item \hfill \textit{10 pts}

An aircraft engine takes in 9.52 kJ of heat energy and discards 6.08 kJ each cycle. Calculate the engine's mechanical work output (in kJ).

\textit{Correct answer:} 3.44

\item \hfill \textit{10 pts}

The figure shows an engine cycle's \emph{pV}-diagram, where \emph{p}$_{1}$ = 93 kPa, \emph{p}$_{2}$ = 177 kPa, \emph{p}$_{3}$ = 283 kPa; \emph{V}$_{1}$ = 1.1 m$^{3}$, and \emph{V}$_{2}$ = 3.7 m$^{3}$. Calculate the net work in kJ done by the engine during one cycle.

 [Image: ME20\_PV\_5.png]

\textit{Correct answer:} 356.2

\item \hfill \textit{10 pts}

A refrigerator transfers a total of 1,305 J of energy into the kitchen while doing 106 J of work. Calculate the coefficient of performance.

\textit{Correct answer:} 11.3113

\item \hfill \textit{10 pts}

A heat pump in heating mode delivers 1,116 J of energy into a home. If the coefficient of performance is 10.9, how many joules of energy does it extract from the cold reservoir?

\textit{Correct answer:} 1,013.6147

\item \hfill \textit{10 pts}

Suppose a heat pump in cooling mode removes 1,050 J of energy from inside the house while ejecting 1,236 J of energy outdoors. Calculate the coefficient of performance.

\textit{Correct answer:} 5.6452

\item \hfill \textit{10 pts}

Suppose an engine's \emph{pV}-diagram takes a rectangular path around pressures 86 kPa and 216 kPa and volumes 0.6 m$^{3}$ and 2.5 m$^{3}$. Calculate the engine's efficiency if the gas is monatomic.

\textit{Correct answer:} 0.2161

\item \hfill \textit{10 pts}

A Carnot engine operates between the temperatures of 334 K and 564 K. How much work does the engine perform when absorbing energy 2,313 J from the hot reservoir?

\textit{Correct answer:} 943.2447

\item \hfill \textit{10 pts}

A Carnot engine ejects 1,171 J to the cold reservoir. If that reservoir has a temperature of 301 K and the engine does 339 J of work, calculate the temperature of the hot reservoir in Kelvin.

\textit{Correct answer:} 388.1383

\item \hfill \textit{10 pts}

Suppose 2.2 kg of a solid substance is at its melting point of -24°C. If the latent heat of fusion is 191,022 J/kg, calculate the change of the substance's entropy when it melts.



Δ\emph{S} = \rule{2.5cm}{0.4pt}J/K

\textit{Correct answer:} 1,687.7446

\item \hfill \textit{10 pts}

Gallium melts to a liquid at 29.76°C. Suppose 39 kg of liquid gallium at this temperature is placed in large room at 18°C. Calculate the change in entropy of the universe during the time that the gallium freezes solid, with neither gallium nor the room significantly changing temperature. \emph{L}$_{f,gallium}$ = 8.04 x 10$^{4}$ J/kg.

\textit{Correct answer:} 418.5395

\item \hfill \textit{10 pts}

Suppose 12 moles of an ideal monatomic gas at 19°C is warmed at constant volume by 42 C°. Calculate the change in entropy of the gas.

\textit{Correct answer:} 20.1125

\item \hfill \textit{10 pts}

An ideal  monatomic gas occupying 1.6 m$^{3}$ is allowed to freely expand into an additional, evacuated volume of 3.9 m$^{3}$. Calculate the change in entropy, if there are 43.2 g of Helium gas, which has molar atomic mass of 4.00 g/mol.

\textit{Correct answer:} 110.8759

\item \hfill \textit{10 pts}

A 5.4-kg block of ice at 0°C melts to water at 0°C and then warms to room temperature at 24°C. Calculate the change in entropy of the universe if the room effectively remains at the same temperature.

[\emph{L}$_{v,water}$ = 2.26 × 10$^{6}$ J/kg, \emph{L}$_{f,water}$ = 3.35 × 10$^{5}$ J/kg, \emph{c}$_{water}$ = 4186 J/(kg⋅K)]

\textit{Correct answer:} 613.5009

\item \hfill \textit{10 pts}

Identify the following statements as true or false.

(A) The entropy of any region in the universe is always increasing.

    [ Select ]

  ["False", "True"]

(B) For any process, the change in the entropy of the universe as a whole is always greater than or equal to zero.

    [ Select ]

  ["False", "True"]

(C) The entropy of the universe always increases during the free expansion of a gas.

    [ Select ]

  ["False", "True"]

(D) A system has zero entropy change during an adiabatic non-isothermal process.

    [ Select ]

  ["True", "False"]

\begin{itemize}
  \item False
  \item True
  \item False
  \item True
  \item False
  \item True
  \item True
  \item False
\end{itemize}

\item \hfill \textit{10 pts}

In a diesel cycle, the thermal process just before the power stroke is
, whereas the thermal process in the Otto cycle is
.



\textbf{Word list: adiabatic, isothermal, isobaric, isochoric, generic}

\begin{itemize}
  \item isobaric
  \item isochoric
\end{itemize}

\end{enumerate}

\end{document}
